\clearpage
\begingroup
\fancyhead[L]{\small REMARKS}
\renewcommand{\headrulewidth}{0pt}
\begingroup
\newcommand{\BVIJ}{\text{\mdseries\AstroSymbols ♃}}
\newcommand{\BVIP}{\text{\mdseries\AstroSymbols ⚴}}
\newcommand{\BVIsin}{\mathop{\mathit{sin}}}
\par\noindent\hspace*{\parindent}The entire calculation was actually completed in July 1813. See the letter to
\textsc{Olbers}, p. 422.

\par\noindent\hspace*{\parindent}The sets of elements printed on p. 529, and the calculation of the
relative positions of the orbits, are in notebook Bc, as is the »preparation
of the positions of Jupiter and Pallas«, that is, the calculation, also mentioned on p. 529,
of the coordinates \(r\), \(v\); \(r'\), \(w'\), \(\beta'\), \(v'\).
The rest of the calculation is carried out in full on a large number of separate sheets,
from which the numbers in the tables on pp. 530--541 and those in the
worked examples for the argument \(5M'-2M\), pp. 542, 545, 548, 550, 552, were taken.
The final results for the perturbations on pp. 543--545,
546--547, 548--549, 550--552 (only \(\delta\log a\)), 553--554, 555--556
are again in the notebook just mentioned. Only the coefficients of
\(\delta\int n\,\mathrm{d}t\) had to be taken from the sheets.
Unfortunately, some terms in the expression for \(\delta\tilde\omega\), p. 547,
were inadvertently omitted. They are therefore given here:
\[
\begin{array}{r@{\ }r@{\ }r}
+\mathit{1\overset{\prime\prime}{,}24}\,\BVIsin(\mathit{14}\BVIJ-\mathit{5}\BVIP+\mathit{253^0}&\mathit{0'}&\mathit{39''})\\
+\mathit{0{,}24}\,\BVIsin(\mathit{14}\BVIJ-\mathit{4}\BVIP+\mathit{251}&\mathit{4}&\mathit{55})\\
+\mathit{0{,}13}\,\BVIsin(\mathit{15}\BVIJ-\mathit{10}\BVIP+\mathit{116}&\mathit{58}&\mathit{57})\\
+\mathit{0{,}19}\,\BVIsin(\mathit{15}\BVIJ-\mathit{9}\BVIP+\mathit{106}&\mathit{32}&\mathit{59})\\
+\mathit{0{,}29}\,\BVIsin(\mathit{15}\BVIJ-\mathit{8}\BVIP+\mathit{99}&\mathit{41}&\mathit{37})\\
+\mathit{0{,}47}\,\BVIsin(\mathit{15}\BVIJ-\mathit{7}\BVIP+\mathit{95}&\mathit{9}&\mathit{38})\\
+\mathit{2{,}59}\,\BVIsin(\mathit{15}\BVIJ-\mathit{6}\BVIP+\mathit{91}&\mathit{52}&\mathit{25})\\
+\mathit{0{,}35}\,\BVIsin(\mathit{15}\BVIJ-\mathit{5}\BVIP+\mathit{268}&\mathit{44}&\mathit{35})\\
+\mathit{0{,}10}\,\sin(\mathit{16}\BVIJ-\mathit{10}\BVIP+\mathit{125}&\mathit{50}&\mathit{29})\\
+\mathit{0{,}15}\,\BVIsin(\mathit{16}\BVIJ-\mathit{9}\BVIP+\mathit{117}&\mathit{59}&\mathit{33})\\
+\mathit{0{,}22}\,\BVIsin(\mathit{16}\BVIJ-\mathit{8}\BVIP+\mathit{109}&\mathit{38}&\mathit{55})\\
+\mathit{0{,}42}\,\BVIsin(\mathit{16}\BVIJ-\mathit{7}\BVIP+\mathit{107}&\mathit{50}&\mathit{23})
\end{array}
\]
\[
\begin{array}{r@{\ }r@{\ }r}
+\mathit{1\overset{\prime\prime}{,}10}\,\BVIsin(\mathit{16}\BVIJ-\mathit{6}\BVIP+\mathit{283^0}&\mathit{10'}&\mathit{1''})\\
+\mathit{0{,}13}\,\BVIsin(\mathit{16}\BVIJ-\mathit{5}\BVIP+\mathit{281}&\mathit{17}&\mathit{22})\\
+\mathit{0{,}12}\,\BVIsin(\mathit{17}\BVIJ-\mathit{9}\BVIP+\mathit{129}&\mathit{36}&\mathit{2})\\
+\mathit{0{,}20}\,\BVIsin(\mathit{17}\BVIJ-\mathit{8}\BVIP+\mathit{125}&\mathit{14}&\mathit{12})\\
+\mathit{0{,}50}\,\BVIsin(\mathit{17}\BVIJ-\mathit{7}\BVIP+\mathit{120}&\mathit{22}&\mathit{4})\\
+\mathit{0{,}22}\,\BVIsin(\mathit{17}\BVIJ-\mathit{6}\BVIP+\mathit{297}&\mathit{40}&\mathit{56})\\
+\mathit{0{,}14}\,\BVIsin(\mathit{18}\BVIJ-\mathit{8}\BVIP+\mathit{137}&\mathit{75}&\mathit{16})\\
+\mathit{0{,}14}\,\BVIsin(\mathit{19}\BVIJ-\mathit{8}\BVIP+\mathit{149}&\mathit{16}&\mathit{20})\\
+\mathit{0{,}16}\,\BVIsin(\mathit{19}\BVIJ-\mathit{7}\BVIP+\mathit{146}&\mathit{41}&\mathit{38})\\
+\mathit{0{,}22}\,\BVIsin(\mathit{20}\BVIJ-\mathit{8}\BVIP+\mathit{161}&\mathit{17}&\mathit{24})\\
+\mathit{0{,}16}\,\BVIsin(\mathit{21}\BVIJ-\mathit{8}\BVIP+\mathit{173}&\mathit{18}&\mathit{28}).
\end{array}
\]

\par\noindent\hspace*{\parindent}The numerical values in note [9.], p. 557, are on the separate
sheets among the integrations for the other equations. Some of the values
had to be supplied.

\par\noindent\hspace*{\parindent}Note [10.] deserves particular interest. It was assembled from
a few very brief notes on two unassuming slips of paper,
one of which
\SourceLeaf{614}comes from a bundle of loose slips said to have been found in \textsc{Gauss}'s wastepaper basket after his death.
The appearance of this slip lends some probability to that account of its origin.
Judging from its contents, however, one can hardly suppose that it was thrown into the basket
shortly before \textsc{Gauss}'s death. In the repeated examinations
of the posthumous papers, it was naturally difficult to avoid some disturbance
of the original order of the individual slips that had been found.
It is scarcely possible to establish with certainty when the two slips were written.
According to the »Journal of calculations on the perturbations of Pallas« printed above,
\textsc{Gauss} could only have learned between 20 July
and 10 August 1812 the numerical value of the coefficient
\(\alpha\) (p. 558), hence after writing the letter to
\textsc{Bessel}, p. 422, and before composing the passage in the »Exposition«, p. 472.
One circumstance that perhaps favours a later date for both slips is that
\textsc{Gauss} made a mistake when taking the numerical values from his
earlier calculations, assigning the constant \(A\) an incorrect value,
namely
\[
A=198^0\,39'\,8'',
\]
\noindent with which he finds
\[
\begin{aligned}
U&=324^0\,15'\,57''\\
\log C&=4{,}26300
\end{aligned}
\]
\noindent and for the period
\[
1026{,}17\ \text{revolutions of Pallas}=399{,}07\ \text{revolutions of Jupiter}.
\]

\par\noindent\hspace*{\parindent}The limits between which \(n\) and \(u\) oscillate (p. 559) were left uncalculated by
\textsc{Gauss}. The value of \(A\) given above would yield:
\[
\begin{array}{ccccc@{\ }l}
\text{for}&\text{the}&\text{limits}&\text{of}&n:&\dfrac{18}{7}n'\pm0\overset{\prime\prime}{,}07219\\[3pt]
\text{»}&\text{»}&\text{»}&\text{»}&u:&\pm37^0\,35'.
\end{array}
\]
\noindent The corrected values have been substituted in the printed text. The difference
is quite considerable. It must be borne in mind, however, that the whole result is very
uncertain, because the value of \(\mu\), and hence that of \(C\), is very difficult to determine with confidence.

\par\noindent\hspace*{\parindent}The originals contain only the equations on p. 558, lines 4,
6, 14 from the top and 8, 6, 5, 4, 3 from the bottom, and p. 559, line 5, as well as the numerical
values of \(\alpha\), \(A\), \(18\BVIJ-7\BVIP\), \(U\), \(14\alpha\), \(\mu\)
and the note »\textit{1026,17 revolutions of Pallas, 399,07 revolutions of Jupiter}«.

\par\noindent\hspace*{\parindent}The investigations on calculating the arguments with allowance for the
great equation, note [11.], are mostly on separate sheets,
the rest in notebook Bc. Only the expressions \(\mathrm{f}(u)\)
and \(\mathrm{F}(u)\), p. 559, were taken from the sheets, and only relations 1) and 3)
from the notebook. Here too the reconstruction had to be fairly free, as the
various surviving calculations give no clear indication of which method
\textsc{Gauss} ultimately regarded as most expedient. The detailed
auxiliary table mentioned on p. 560 is in notebook Bc. It seemed
unnecessary to print it. The numerical values for comparison with the oppositions, p. 560,
are likewise scattered through notebook Bc. The amounts of the perturbations here are calculated from the
Manipel mentioned above, with \textsc{Gauss} receiving assistance from \textsc{Wachter}.
The surviving calculations cover only the four elements treated on p. 560.

\par\noindent\hspace*{\parindent}Note [12.], which on p. 561 is inadvertently labelled [11.],
contains the comparison with the first nine oppositions, which \textsc{Gauss} mentions in
his letter to \textsc{Bessel} of 21 March 1843 (p. 438). It too is carried out on
separate sheets from which only the numbers could be taken. On those
sheets \textsc{Gauss} later added a few more numbers relating to
the tenth opposition. In addition to
\SourceLeaf{615}those printed here, there are, as mentioned at the beginning of these remarks,
very many further comparisons, also with later oppositions.
\par
\endgroup
\clearpage
\endgroup
